% !TEX program = xelatex
% Báo cáo kỹ thuật: AMP/GAIL định hướng policy đi bộ G1 trên mjlab.
% Biên dịch: cd report && latexmk -xelatex main.tex
\documentclass[11pt,a4paper]{article}

\usepackage{fontspec}
\usepackage{babel}
\babelprovide[import,main]{vietnamese}
% Font của TeX Live (có sẵn trên Overleaf), nạp theo tên file.
\setmainfont{texgyretermes}[Extension=.otf, UprightFont=*-regular,
  BoldFont=*-bold, ItalicFont=*-italic, BoldItalicFont=*-bolditalic]
\setsansfont{texgyreheros}[Extension=.otf, UprightFont=*-regular,
  BoldFont=*-bold, ItalicFont=*-italic, BoldItalicFont=*-bolditalic]
\setmonofont{lmmono10}[Extension=.otf, UprightFont=*-regular, Scale=0.95]

\usepackage[margin=2.3cm]{geometry}
\usepackage{amsmath,amssymb,mathtools}
\usepackage{booktabs,multirow,array}
\usepackage{graphicx}
\usepackage[table]{xcolor}
\usepackage{float}
\usepackage{caption,subcaption}
\usepackage{enumitem}
\usepackage{listings}
\usepackage[most]{tcolorbox}
\usepackage{tikz}
\usetikzlibrary{arrows.meta,positioning,fit,calc}
\usepackage[colorlinks=true,linkcolor=blue!50!black,citecolor=green!40!black,
  urlcolor=blue!60!black]{hyperref}
\usepackage{cleveref}
\crefname{figure}{Hình}{Hình}\Crefname{figure}{Hình}{Hình}
\crefname{table}{Bảng}{Bảng}\Crefname{table}{Bảng}{Bảng}
\crefname{section}{Mục}{Mục}\Crefname{section}{Mục}{Mục}
\crefname{equation}{}{}\Crefname{equation}{}{}

% ---------------------------------------------------------------------------
% Macro
% ---------------------------------------------------------------------------
\newcommand{\todo}[1]{\textcolor{red}{\textbf{[TODO:} #1\textbf{]}}}
\newcommand{\E}{\mathbb{E}}
\newcommand{\R}{\mathbb{R}}
\newcommand{\pit}{\pi_\theta}
\newcommand{\code}[1]{\texttt{#1}}
\DeclareMathOperator*{\argmax}{arg\,max}
\DeclareMathOperator*{\argmin}{arg\,min}

\lstset{basicstyle=\ttfamily\small,breaklines=true,frame=single,
  backgroundcolor=\color{gray!6},columns=fullflexible}

\newtcolorbox{algobox}[1]{colback=gray!4,colframe=gray!60,
  title={#1},fonttitle=\bfseries,left=4pt,right=4pt}

\newtcolorbox{notebox}{colback=blue!3,colframe=blue!40,left=4pt,right=4pt}

% Chèn hình nếu đã có, ngược lại hiện khung giữ chỗ.
\newcommand{\maybefig}[2]{%
  \IfFileExists{figures/#1}{\includegraphics[width=#2]{figures/#1}}{%
    \fbox{\parbox[c][4.5cm][c]{\dimexpr#2-2\fboxsep-2\fboxrule}{\centering\todo{hình \texttt{figures/\detokenize{#1}}}}}}}

% Một dải ảnh chụp robot: nhãn + ảnh.
\newcommand{\snap}[2]{\makebox[2.2cm][l]{\small #1}%
  \includegraphics[width=0.8\linewidth]{figures/snap_#2_crop.png}\\[1pt]}

% Hộp mô tả một bước: tiêu đề, (vào / làm / ra), giải thích công thức.
\newcommand{\io}[3]{%
\begin{tcolorbox}[colback=gray!3,colframe=gray!50,left=4pt,right=4pt,top=2pt,
  bottom=2pt,title={#1},fonttitle=\bfseries\small]
\small #2
\par\smallskip #3
\end{tcolorbox}}

\title{\textbf{Định hướng policy đi bộ cho robot hình người Unitree G1\\
bằng Adversarial Imitation Learning trên mjlab}\\[4pt]
\large Báo cáo kỹ thuật -- Sprint 1}
\author{Trần Phú Nghĩa}
\date{25/09/2026}

